\subsection{通过平移分解计算复合积}

设

\[
0 \rightarrow \mathrm{M} \xrightarrow {f} \mathrm{K} _ {n} \rightarrow \mathrm{K} _ {n - 1} \rightarrow \dots \rightarrow \mathrm{K} _ {1} \xrightarrow {g} \mathrm{M} ^ {\prime} \rightarrow 0\tag{16}
\]

是左 A-模的正合序列，且 \(\theta\in\mathrm{Ext}_{\mathbf{A}}^{n}(\mathbf{M}^{\prime},\mathbf{M})\) 是相应的类。

设 \(a:(\mathbf{R},d)\to \mathbf{M}\) 是 \(\mathbf{M}\) 的左分解；因此我们有一个正合序列

\[
\rightarrow \mathrm{R} _ {k} \xrightarrow {d _ {k}} \mathrm{R} _ {k - 1} \rightarrow \dots \xrightarrow {d _ {1}} \mathrm{R} _ {0} \xrightarrow {a _ {0}} \mathrm{M} \rightarrow 0.
\]

并且通过 n 次平移（X，第 26 页）得到一个正合序列

\[
\rightarrow \mathrm{R} _ {k} \xrightarrow {(- 1) ^ {n} d _ {k}} \mathrm{R} _ {k - 1} \rightarrow \dots \xrightarrow {(- 1) ^ {n} d _ {1}} \mathrm{R} _ {0} \xrightarrow {(- 1) ^ {n} a _ {0}} \mathrm{M} \rightarrow 0.\tag{17}
\]

由 (16) 和 (17) 我们推出一个正合序列

\[
\rightarrow \mathrm{R} _ {k} \xrightarrow {(- 1) ^ {n} d _ {k}} \mathrm{R} _ {k - 1} \rightarrow \dots \xrightarrow {(- 1) ^ {n} d _ {1}} \mathrm{R} _ {0} \xrightarrow {(- 1) ^ {n} f \circ a _ {0}} \mathrm{K} _ {n} \rightarrow \mathrm{K} _ {n - 1} \rightarrow \dots \rightarrow \mathrm{K} _ {1} \rightarrow \mathrm{M} ^ {\prime} \rightarrow 0
\]

由此得到一个分解 \(\mathbf{R}'\) ，即 \(\mathbf{M}'\) 的分解；记 \(\varphi: \mathbf{R}' \to \mathbf{R}(-n)\) 为由 \(\varphi_k = 1_{\mathbf{R}_{k-n}}\) （当 \(k \geqslant n\) ）所确定的态射 .

若 N 是左 A-模且 P 是右 A-模，则我们有同态

\[
\mathrm{H} \left(1 _ {\mathrm{P}} \otimes \varphi\right): \mathrm{H} \left(\mathrm{P} \otimes_ {\mathrm{A}} \mathrm{R} ^ {\prime}\right)\rightarrow \mathrm{H} \left(\mathrm{P} \otimes_ {\mathrm{A}} \mathrm{R}\right) (- n)
\]

\[
\mathrm{H} \left(\operatorname{Homgr} _ {\mathbf {A}} (\varphi , 1 _ {\mathbf {N}})\right): \mathrm{H} \left(\operatorname{Homgr} _ {\mathbf {A}} (\mathrm{R}, \mathrm{N})\right) (n) \rightarrow \mathrm{H} \left(\operatorname{Homgr} _ {\mathbf {A}} \left(\mathrm{R} ^ {\prime}, \mathrm{N}\right)\right).
\]

设 k 为整数。

命题 8. --- a) 下图，其中 \(h_{\theta}(\alpha) = \theta \circ \alpha\) ，是交换的

\[
\begin{array}{c} \operatorname{Tor} _ {k + n} ^ {\mathrm{A}} (\mathrm{P}, \mathrm{M} ^ {\prime}) \xrightarrow {h _ {\theta}} \operatorname{Tor} _ {k} ^ {\mathrm{A}} (\mathrm{P}, \mathrm{M}) \\ \psi_ {k + n} (\mathrm{P}, \mathrm{R} ^ {\prime}) \Big | _ {\downarrow} \quad \Big | _ {\downarrow} \psi_ {k} (\mathrm{P}, \mathrm{R}) \\ \mathrm{H} _ {k + n} (\mathrm{P} \otimes_ {\mathrm{A}} \mathrm{R} ^ {\prime}) \xrightarrow {\mathrm{H} _ {k + n} (1 \otimes \varphi)} \mathrm{H} _ {k} (\mathrm{P} \otimes_ {\mathrm{A}} \mathrm{R}) \end{array}
\]

\begin{enumerate}
\def\labelenumi{\alph{enumi})}
\setcounter{enumi}{1}
\tightlist
\item
  下图，其中 \(\delta_{\theta}(\beta) = \beta \circ \theta\) ，是交换的
\end{enumerate}

\[
\begin{array}{c} \mathrm{H} ^ {k} (\text {Homgr} _ {\mathrm{A}} (\mathrm{R}, \mathrm{N})) \xrightarrow {\mathrm{H} ^ {k + n} (\text {Homgr} _ {\mathrm{A}} (\varphi , 1))} \mathrm{H} ^ {k + n} (\text {Homgr} _ {\mathrm{A}} (\mathrm{R} ^ {\prime}, \mathrm{N})) \\ \varphi^ {k} (\mathrm{R}, \mathrm{N}) \Biggl \downarrow \\ \operatorname{Ext} _ {\mathrm{A}} ^ {k} (\mathrm{M}, \mathrm{N}) \xrightarrow {\delta_ {\theta}} \operatorname{Ext} _ {\mathrm{A}} ^ {k + n} (\mathrm{M} ^ {\prime}, \mathrm{N}). \end{array}
\]

设 \(\alpha : L(M) \to R\) 是复形的态射，使得 \(a \circ \alpha = p_{M}\) 并设

\[
\begin{array}{c} \mathrm{L} _ {n} (\mathrm{M} ^ {\prime}) \longrightarrow \mathrm{L} _ {n - 1} (\mathrm{M} ^ {\prime}) \longrightarrow \dots \longrightarrow \mathrm{L} _ {0} (\mathrm{M} ^ {\prime}) \longrightarrow \mathrm{M} ^ {\prime} \longrightarrow 0 \\ u _ {n} \Big \downarrow \qquad \qquad \qquad \Big \downarrow u _ {n - 1} \qquad \qquad \qquad \qquad \qquad \qquad \qquad \qquad \qquad \qquad \qquad \qquad \qquad \qquad \qquad \qquad \qquad \qquad \qquad \qquad \qquad \qquad \qquad \qquad \qquad \qquad \qquad \qquad \qquad \qquad \qquad \qquad \qquad \qquad \\ 0 \longrightarrow \mathrm{M} \xrightarrow {f} \mathrm{K} _ {n} \longrightarrow \dots \longrightarrow \mathrm{K} _ {1} \longrightarrow \mathrm{M} ^ {\prime} \longrightarrow 0 \end{array}
\]

一个交换图；让我们选取一个同态 \(v_{n}: \mathrm{L}_{n}(\mathrm{M}^{\prime}) \to \mathrm{L}_{0}(\mathrm{M})\) 使得 \(p_{\mathbf{M}} \circ v_{n} = (-1)^{n} u_{n}\) ；由 X，第 47 页，命题 1, a)， \(v_{n}\) 可延拓为复形的态射 \(v: \mathrm{L}(\mathrm{M}^{\prime}) \to \mathrm{L}(\mathrm{M})(-n)\) ，且 \(\theta\) 是典范同构 \(\mathrm{H}^{n}(\mathrm{Homgr}_{\mathbf{A}}(\mathrm{L}(\mathrm{M}^{\prime}), \mathrm{L}(\mathrm{M}))) \to \mathrm{Ext}_{\mathbf{A}}^{n}(\mathbf{M}^{\prime}, \mathbf{M})\) 下 \(v\) 的类的像 (X，第 117 页，注记 1)。定义复形的态射 \(\beta: \mathrm{L}(\mathrm{M}^{\prime}) \to \mathrm{R}^{\prime}\) 为： \(\beta_{p} = u_{p}\) （当 \(p \leqslant n - 1\)）， \(\beta_{p} = \alpha_{p-n} \circ v_{p}\) （当 \(p \geqslant n\)），并且我们有

\[
\varphi \circ \beta = \alpha (- n) \circ v.
\]

另一方面，根据映射 \(\varphi\) 与 \(\psi\) 的定义：

\[
\begin{array}{r l r} & {\psi_ {k} (\mathrm{P}, \mathrm{R}) = \mathrm{H} _ {k} (p _ {\mathrm{P}} \otimes \alpha),} & {\varphi^ {k} (\mathrm{R}, \mathrm{N}) = \mathrm{H} ^ {k} (\mathrm{Homgr} _ {\mathrm{A}} (\alpha , e _ {\mathrm{N}})),} \\ & {\psi_ {k + n} (\mathrm{P}, \mathrm{R} ^ {\prime}) = \mathrm{H} _ {k + n} (p _ {\mathrm{P}} \otimes \beta),} & {\varphi^ {k + n} (\mathrm{R} ^ {\prime}, \tilde {\mathrm{N}}) = \mathrm{H} ^ {k + n} (\mathrm{Homgr} _ {\mathrm{A}} (\beta , e _ {\mathrm{N}})).} \end{array}
\]

最后，根据复合乘积的定义，有

\[
h _ {\theta} = \mathrm{H} (1 _ {L (\mathbb {P})} \otimes v), \quad \delta_ {\theta} = \mathrm{H} (\operatorname{Homgr} _ {\mathbf {A}} (v, 1 _ {I (\mathbb {N})}))
\]

因此，有以下等式

\[
\begin{array}{r l} \psi_ {k} (\mathrm{P}, \mathrm{R}) \circ h _ {\theta} & = \mathrm{H} _ {k} (p _ {\mathrm{P}} \otimes \alpha) \circ \mathrm{H} _ {k} (1 _ {\mathrm{L(P)}} \otimes v) = \mathrm{H} _ {k} (p _ {\mathrm{P}} \otimes (\alpha \circ v)) = \mathrm{H} _ {k + n} (p _ {\mathrm{P}} \otimes (\varphi \circ \beta)) \\ & = \mathrm{H} _ {k + n} (1 \otimes \varphi) \circ \mathrm{H} _ {k + n} (p _ {\mathrm{P}} \circ \beta) = \mathrm{H} _ {k + n} (1 \otimes \varphi) \circ \psi_ {k + n} (\mathrm{P}, \mathrm{R} ^ {\prime}), \end{array}
\]

由此得 a)；b) 的证明类似。

注记. --- 由交换同构，我们从 \(a\) ) 推出关于右 A-模的正合序列 (16) 情形的类似结论。

现在设 \(b: \mathbf{M}' \to \mathbf{E}'\) 为 \(\mathbf{M}'\) 的右分解；因此我们有一个正合序列

\[
0 \rightarrow \mathbf {M} ^ {\prime} \xrightarrow {\dot {b} ^ {0}} \mathrm{E} ^ {\prime 0} \xrightarrow {\delta^ {0}} \mathrm{E} ^ {\prime 1} \rightarrow \dots \rightarrow \mathrm{E} ^ {\prime k} \xrightarrow {\delta^ {k}} \mathrm{E} ^ {\prime k + 1}
\]

由此得到一个正合序列

\[
0 \rightarrow \mathrm{M} \xrightarrow {f} \mathrm{K} _ {n} \rightarrow \mathrm{K} _ {n - 1} \rightarrow \dots \rightarrow \mathrm{K} _ {1} \xrightarrow {(- 1) ^ {n} b ^ {0} \circ g} \mathrm{E} ^ {\prime 0} \xrightarrow {(- 1) ^ {n} \delta^ {0}} \mathrm{E} ^ {\prime 1} \rightarrow \dots
\]

对应于 M 的一个右分解 E；记 \(\sigma : E'(n) \to E\) 为由 \(\sigma^k = 1_{E'k - n}\) （当 \(k \geqslant n\) ）所确定的态射。因此我们有同态

\[
\mathrm{H} \left(\operatorname{Homgr} _ {\mathrm{A}} \left(1 _ {\mathrm{N}}, \sigma\right)\right): \mathrm{H} \left(\operatorname{Homgr} _ {\mathrm{A}} \left(\mathrm{N}, \mathrm{E} ^ {\prime}\right)\right) (n) \rightarrow \mathrm{H} \left(\operatorname{Homgr} _ {\mathrm{A}} (\mathrm{N}, \mathrm{E})\right).
\]

命题 9. --- 下图，其中 \(\gamma_{\theta}(\alpha) = \theta \circ \alpha\) ，是交换的：

\[
\begin{array}{c} \mathrm{H} ^ {k} (\text {Homgr} _ {\mathrm{A}} (\mathrm{N}, \mathrm{E} ^ {\prime})) \xrightarrow {\mathrm{H} ^ {k + n} (\text {Homgr} _ {\mathrm{A}} (1 _ {\mathrm{N}} , \sigma))} \mathrm{H} ^ {k + n} (\text {Homgr} _ {\mathrm{A}} (\mathrm{N}, \mathrm{E})) \\ \varphi^ {k} (\mathrm{N}, \mathrm{E} ^ {\prime}) \Big \downarrow \\ \operatorname{Ext} _ {\mathrm{A}} ^ {k} (\mathrm{N}, \mathrm{M} ^ {\prime}) \xrightarrow {\gamma_ {\theta}} \operatorname{Ext} _ {\mathrm{A}} ^ {k + n} (\mathrm{N}, \mathrm{M}). \end{array}
\]

此命题的证明方式与命题 8 类似。

